\begin{tikzpicture}[
  font=\scriptsize,
  panah/.style={-{Stealth[length=2mm]}, thick},
  kotak/.style={draw, rounded corners, fill=blue!8, minimum width=2.1cm,
    minimum height=0.9cm, align=flush center, text width=1.95cm},
  basisdata/.style={draw, cylinder, shape border rotate=90, aspect=0.22,
    minimum width=1.9cm, minimum height=1.15cm, align=flush center, text width=1.7cm},
  judul/.style={font=\tiny\bfseries, gray!60!black, align=flush center}
]
  % Sisi perintah: layanan operasional menulis data miliknya sendiri.
  \node[judul] at (0,2.75) {Sisi perintah (\textit{command})};
  \node[kotak] (kunjungan) at (0,1.4) {Layanan\\Kunjungan};
  \node[kotak] (tagihan) at (0,0) {Layanan\\Tagihan};
  \node[kotak] (lab) at (0,-1.4) {Layanan\\Laboratorium};

  % Event mengalir ke broker, lalu diproyeksikan menjadi read model.
  \node[kotak, minimum height=3.7cm, minimum width=1.6cm, text width=1.4cm]
    (broker) at (2.85,0) {\textit{Event broker}};
  \foreach \asal in {kunjungan, tagihan, lab}
    \draw[panah] (\asal.east) -- (\asal.east -| broker.west);
  \node[kotak, fill=orange!12] (proyeksi) at (5.4,0)
    {Pembangun \textit{read model}\\{\tiny\mdseries proyeksi \textit{event}}};
  \draw[panah] (broker.east |- proyeksi) -- (proyeksi.west);

  % Sisi baca: read model melayani laporan tanpa menyentuh layanan operasional.
  \node[judul] at (8.2,2.75) {Sisi baca (\textit{query})};
  \node[basisdata, fill=green!10] (readmodel) at (8.2,0)
    {\textit{Read model}\\laporan};
  \node[kotak, fill=praditagreen!15] (dashboard) at (8.2,1.95)
    {\textit{Dashboard}\\manajemen};
  \draw[panah] (proyeksi.east) -- (readmodel.west);
  \draw[panah] (dashboard.south) -- node[right, font=\tiny] {\textit{query}}
    (readmodel.north);

  % Batasan kesegaran ditulis di bawah gambar.
  \node[align=flush center, font=\tiny\mdseries] at (4.1,-2.55)
    {\textit{Read model} tertinggal beberapa detik sampai beberapa menit dari data
     operasional,\\dan dapat dibangun ulang dari \textit{event} bila rusak};
\end{tikzpicture}
